\begin{textsample}{POS Dim 1 – 2013\_c6 – Score -1.00 – 2013\_c6/enade\_2013\_tec\_agro\_5e89c31973.txt}  \label{ex:f1_pos_176}
CAPÍTULO 3 ANÁLISE TÉCNICA DA PROVA Este capítulo tem por objetivo apresentar o desempenho dos estudantes concluintes de Tecnologia em Agronegócio no ENADE/2013.
Para isso, foram calculadas as estatísticas básicas da prova em seu todo ( seção 3.1.1 ), bem como as estatísticas dos componentes relacionadas ao Componente de Formação Geral ( seção 3.1.2 ) e ao de Conhecimento Específico da Área ( seção 3.1.3 ).
Dadas as suas características, foram analisadas em separado, as questões objetivas ( seção 3.2 ) e as questões discursivas ( seção 3.3 ).
Para as questões objetivas, foram disponibilizados os índices de facilidade e o de discriminação bisserial, também em separado para os Componentes de Formação Geral ( seção 3.2.1 ) e de Conhecimento Específico ( seção 3.2.2 ).
De cada componente, uma das questões foi escolhida para exemplificar a análise gráfica relacionando as alternativas escolhidas pelos estudantes ( inclusive o gabarito ) com o número de acertos no componente.
O Anexo I apresenta a íntegra da análise gráfica para todas as questões objetivas.
Para cada uma das questões discursivas, os conteúdos dos tipos mais comuns de respostas dos alunos são apresentados e comparados com o padrão de resposta esperado ( ver Anexo VIII com o padrão de respostas ).
Tomando como base as duas questões discursivas do Componente de Formação Geral, a seção 3.3.1.6 apresenta comentários sobre a correção das respostas com respeito à Língua Portuguesa.
Nas tabelas, são apresentados o tamanho da população inscrita e de presentes, e as seguintes estatísticas das notas¹¹ : média do desempenho na prova, erro padrão da média, desvio padrão, nota mínima, mediana e nota máxima.
Tais estatísticas contemplam o total de estudantes concluintes da Área de Tecnologia em Agronegócio inscritos e presentes à prova do ENADE/2013, tendo -se em vista as seguintes agregações : ( a ) as Grandes Regiões e o país como um todo ; ( b ) a Categoria Administrativa ; e ( c ) a Organização Acadêmica.
Em relação aos gráficos de distribuição de notas, o intervalo considerado foi de 10 unidades, aberto à esquerda e fechado à direita, com exceção do primeiro intervalo, [ 0 ; 10 ], fechado em ambos os extremos.
Para os gráficos de distribuição das notas das questões discursivas, foram consideradas mais duas categorias : questão em branco e nota zero. ¹¹ Essas estatísticas e outras estão definidas no Capítulo 1. 3.1 ESTATÍSTICAS BÁSICAS DA PROVA 3.1.1 Estatísticas Básicas Gerais A Tabela 3.1 apresenta as estatísticas básicas da prova por grande Região.
A população total de inscritos foi de 1.739.
Destes, 1.452 estiveram presentes, sendo 16,5 \% o índice de não comparecimento.
A Região de maior abstenção foi a Centro-Oeste ( 27,4 \% ) e a de menor abstenção foi a Nordeste ( 5,9 \% ).
A média das notas da prova como um todo ( nas seções seguintes serão analisados os Componentes de Formação Geral e de Conhecimento Específico ) foi 49,6, sendo que os alunos da região Centro-Oeste obtiveram a média mais baixa ( 45,4 ) e os da região Sudeste obtiveram a média mais alta ( 50,8 ).
As demais médias foram : 48,5 na região Norte, 50,0 na região Nordeste e 48,8 na região Sul.
O desvio padrão para o Brasil como um todo foi 12,7, sendo o maior desvio padrão encontrado na região Centro-Oeste ( 13,5 ) e o menor na região Nordeste ( 10,8 ), indicando uma menor dispersão das notas desta última região.
A região que obteve a maior nota máxima foi a Sudeste ( 82,8 ), ao passo que a região que atingiu a menor nota máxima foi a Nordeste ( 73,7 ).
A mediana do Brasil como um todo foi 50,0, sendo a maior mediana obtida na região Sudeste ( 51,0 ) e a menor obtida na Centro-Oeste ( 46,9 ).
A região que obteve a maior nota mínima foi a Nordeste ( 22,4 ), ao passo que a região que atingiu a menor foi a Sudeste ( 5,6 ).
O comportamento das notas dos estudantes de todo o Brasil pode ser observado no Gráfico 3.1 que apresenta um histograma com a distribuição das mesmas.
Essa é uma distribuição unimodal com moda no intervalo ( 40 ; 50 ].
Destaca -se, também, o intervalo ( 50 ; 60 ] com frequência relativa muito próxima à modal.
O coeficiente de assimetria da distribuição das notas é negativo ( -0,30 ) e não podemos considerar que a distribuição é simétrica.
As distribuições por Grande Região também apresentam assimetria negativa, concentração pouco maior do lado direito do histograma e mais espalhada do lado esquerdo.
Os Gráficos 3.2, 3.3 e 3.4 apresentam informações referentes à média da nota final dos participantes, desagregando os resultados de acordo com, respectivamente, as Grandes Regiões do país, a Categoria Administrativa e a Organização Acadêmica.
Os gráficos apresentam o valor da média das notas como uma barra e os extremos do intervalo de confiança de 95 \% como linhas verticais unidas por uma linha horizontal na forma da letra H maiúscula.
Considerando -se o gráfico de notas segundo Grande Região ( Gráfico 3.2 ), observa -se que existe diferença estatisticamente significativa ao nível de 95 \% entre a maior média, obtida na região Sudeste ( 50,8 ) e a menor, obtida na região Centro-Oeste ( 45,4 ).
Levando -se em conta os estudantes por Categorias Administrativas ( Gráfico 3.3 ), observa -se que existe diferença estatisticamente significativas entre as médias das notas das IES Públicas e Privadas.
Pode ser observado ainda que a média dos alunos de IES Públicas ( 52,3 ) é maior que a dos alunos de IES Privadas ( 46,5 ).
A diferença entre as médias das regiões Sudeste e Centro-Oeste ( 5,4 ), a maior e a menor média, é inferior a diferença entre IES Públicas e Privadas ( 5,8 ).
Tendo como foco o Gráfico 3.4, que apresenta as notas médias das provas segundo Organização Acadêmica, constata -se que a diferença da média das notas dos estudantes provenientes de Universidades não é estatisticamente significativa ao nível de 95 \% das médias das notas dos estudantes provenientes de Centros Universitários e Faculdades, assim como a diferença entre a média dos estudantes de Centros Universitários e Faculdades.
A maior média foi obtida pelos estudantes de Faculdades ( 50,3 ), e a menor, pelos de Centros Universitários ( 47,4 ). 3.1.2 Estatísticas Básicas no Componente de Formação Geral A Tabela 3.2 apresenta as estatísticas básicas em relação ao componente da prova que avalia a Formação Geral dos estudantes concluintes.
Os alunos de todo Brasil obtiveram desempenho médio de 47,1.
Quanto à variabilidade, o desvio padrão das notas dos estudantes do Brasil como um todo foi 13,9.
A maior média foi obtida na região Nordeste ( 51,9 ), e a menor, na região Centro-Oeste ( 43,4 ).
As demais médias foram : 47,5 na região Norte, 47,9 na região Sudeste e 45,1 na região Sul.
Já o maior desvio padrão foi obtido na região Centro-Oeste ( 14,6 ) e o menor na região Sul ( 13,1 ).
Os demais desvios padrões foram : 14,2 na região Norte, 13,5 na região Nordeste e 13,9 na região Sudeste.
A maior nota no Componente de Formação Geral da prova do ENADE foi obtida por pelo menos um aluno da região Sudeste ( 83,1 ) enquanto que a menor nota máxima foi obtida na região Nordeste ( 78,5 ).
Nas outras regiões as notas máximas foram : 78,9 tanto para a região Norte quanto para a região Sul e 81,9 na região Centro-Oeste.
A mediana do Brasil como um todo foi 46,8, sendo a menor mediana encontrada na região Centro-Oeste ( 43,0 ) e a maior encontrada na região Nordeste ( 51,3 ).
A menor nota mínima nesta parte foi obtida na região Sudeste ( 0,0 ), enquanto a maior nota mínima foi obtida na região Norte ( 15,0 ).
Nas outras regiões as notas mínimas foram : 13,5 na região Nordeste e 7,5 tanto para a região Sul quanto para a região Centro-Oeste.
O Gráfico 3.5 propicia a avaliação do desempenho dos estudantes no Componente de Formação Geral a partir do histograma da distribuição das notas correspondentes.
A distribuição é unimodal, com moda em ( 40 ; 50 ], assim como na prova como um todo.
Nota -se, ainda, que no Gráfico 3.5 as notas apresentam uma maior dispersão do que no Gráfico 3.1 ( distribuição das notas da prova ), confirmado pela comparação dos desvios padrões : 12,7 para a nota da prova como um todo e 13,9 para o Componente de Formação Geral.
Para o Componente de Formação Geral, o coeficiente de assimetria da distribuição das notas dos estudantes é negativo ( -0,08 ).
Por ser um coeficiente pequeno, a distribuição apresenta uma leve concentração à direita e cauda maior à esquerda, porém podemos considerar que a distribuição é aproximadamente simétrica.
Em quase todas as Grandes Regiões os histogramas também possuem assimetria levemente negativa ( -0,26 para o Nordeste, -0,14 para o Sudeste e -0,02 para o Sul ), exceto nas regiões Norte ( 0,15 ) e Centro-Oeste ( 0,08 ) que possuem assimetria levemente positiva.
Nos Gráficos 3.6, 3.7 e 3.8 são apresentadas as informações referentes ao desempenho dos concluintes no Componente de Formação Geral, em diferentes agregações : Grande Região do país, Categoria Administrativa e Organização Acadêmica.
Observa -se pelo Gráfico 3.6 que existe diferença estatisticamente significativa ao nível de 95 \% entre a maior média das notas na Componente de Formação Geral, obtida na região Nordeste ( 51,9 ) e a menor, obtida na região Centro-Oeste ( 43,4 ).
Vemos que o intervalo de confiança mais largo é o da região Norte ; já o intervalo mais estreito é observado na região Sudeste.
Este fato está relacionado, também, com o tamanho da população envolvida, menor na região Norte do que na Sudeste.
No Gráfico 3.7, que representa as notas médias no Componente de Formação Geral segundo Categoria Administrativa do país, observa -se que existe diferença estatisticamente significativa entre as médias.
Assim como ocorreu para as notas da prova como um todo, para o Componente de Formação Geral, os concluintes de Tecnologia em Agronegócio, das IES Públicas ( 49,0 ) obtiveram uma média maior do que os das IES Privadas ( 44,8 ).
Considerando -se o tipo de Organização Acadêmica, nota -se, no Gráfico 3.8, que existe uma diferença estatisticamente significativa entre a maior e a menor média.
Nas Faculdades, a média ( 48,0 ) é mais elevada do que em Centros Universitários ( 44,1 ) e Universidades ( 46,5 ). 3.1.3 Estatísticas Básicas do Componente de Conhecimento Específico A Tabela 3.3 apresenta as estatísticas básicas referentes ao Componente de Conhecimento Específico da área de Tecnologia em Agronegócio.
A média do desempenho dos alunos do Brasil como um todo foi 50,4.
A maior média foi obtida na região Sudeste ( 51,7 ) e a menor, na região Centro-Oeste ( 46,0 ).
As demais médias foram : 48,9 na região Norte, 49,4 na região Nordeste e 50,1 na região Sul.
Quanto à variabilidade das notas, o desvio padrão do Brasil como um todo foi 14,3, sendo o maior desvio padrão observado na região Norte ( 15,4 ) e o menor na região Nordeste ( 12,4 ).
Os demais desvios padrões foram : 13,9 na região Sudeste, 14,6 na região Sul e 15,2 na região Centro-Oeste.
A mediana das notas dos estudantes de todo o Brasil foi 51,2.
A maior mediana ocorreu na região Sudeste ( 52,7 ) e a menor na região Centro-Oeste ( 48,2 ).
As demais medianas foram 49,2 na região Norte, 49,0 na região Nordeste e 50,1 na região Sul.
A nota máxima do Brasil como um todo foi 87,7, sendo obtida por pelo menos um aluno da região Sudeste.
As demais notas máximas foram : 83,8 na região Norte, 74,3 na região Nordeste, 81,1 na região Sul e 83,6 na região Centro-Oeste.
A menor nota mínima foi obtida na região Sudeste ( 0,0 ), enquanto a maior nota mínima foi obtida na região Norte ( 17,9 ).
As demais notas mínimas foram : 16,9 na região Nordeste, 4,5 na região Sul e 10,2 na região Centro-Oeste.
Assim como os Gráficos 3.1 e 3.5, o Gráfico 3.9, apresentado a seguir, proporciona uma avaliação do desempenho de concluintes em relação ao Componente de Conhecimento Específico com um histograma da distribuição das notas correspondentes.
Esta também é uma distribuição unimodal, e o grupo modal é o ( 40 ; 50 ], o mesmo grupo modal da Parte de Formação Geral.
O coeficiente de assimetria da distribuição das notas do Componente de Conhecimento Específico é negativo ( -0,33 ).
Nota -se pelo histograma ( Gráfico 3.9 ) que esta distribuição apresenta uma leve concentração à direita e cauda maior à esquerda.
Em todas as Grandes Regiões os histogramas também possuem assimetria negativa ( entre -0,37 para o Sudeste e -0,06 para o Norte ).
Os Gráficos 3.10, 3.11 e 3.12 apresentam, respectivamente, uma comparação dos resultados em relação à Grande Região do país, à Categoria Administrativa e à Organização Acadêmica, agora levando em conta o desempenho de alunos no Componente de Conhecimento Específico da prova.
Pelo Gráfico 3.10, observa -se que existe diferença estatisticamente significativa entre a maior e a menor média das notas no Componente de Conhecimento Específico, obtidas nas regiões Sudeste ( 51,7 ) e Centro-Oeste ( 46,0 ) respectivamente.
Quanto à Categoria Administrativa ( Gráfico 3.11 ), observa -se um comportamento semelhante àquele da Parte de Formação Geral e à prova como um todo, ou seja, existe diferença estatisticamente significativa entre as médias das IES Públicas e Privadas, sendo que a maior média foi obtida por alunos de IES Públicas de ensino ( 53,4 ).
Quanto ao Gráfico 3.12, observa -se, que, assim como na prova como um todo, não existe diferença estatisticamente significativa ao nível de 95 \% entre as notas no Componente de Conhecimento Específico dos diferentes tipos de Organização Acadêmica.
Sendo que a média dos concluintes das Faculdades ( 51,0 ) foi maior do que de Centros Universitários ( 48,5 ) e de Universidades ( 50,1 ). 3.2 ANÁLISE DAS QUESTÕES OBJETIVAS 3.2.1 Componente de Formação Geral A Tabela 3.4 apresenta as estatísticas básicas relativas às oito questões objetivas do componente da prova que abrange a Formação Geral dos estudantes.
A média do Brasil foi 49,0.
A menor média foi encontrada na região Centro-Oeste ( 45,2 ) e a maior na região Sudeste ( 50,5 ).
As demais médias foram 47,8 na região Norte, 50,4 na região Nordeste e 47,1 na região Sul.
O desvio padrão do Brasil foi 17,9, sendo o maior desvio padrão encontrado na região Centro-Oeste ( 19,0 ) e o menor na região Sul ( 16,9 ).
Os demais desvios foram : 17,6 na região Norte, 17,3 na região Nordeste e 18,0 na região Sudeste.
A mediana das notas dos estudantes foi 50,0 em todas as regiões do Brasil, sem exceção.
A nota máxima do Brasil como um todo foi 100,0, sendo obtida por pelo menos um aluno da região Sudeste, enquanto nas demais regiões do Brasil as notas máximas foram 87,5.
A menor nota mínima foi obtida na região Sudeste ( 0,0 ), enquanto a maior nota mínima foi obtida na região Norte ( 25,0 ).
Todas as demais regiões, Nordeste, Sul e Centro-Oeste, obtiveram a mesma nota mínima : 12,5.
A Tabela 3.5 apresenta o índice de facilidade e o índice de discriminação ( ponto bisserial ) para cada uma das questões objetivas do Componente de Formação Geral.
Quanto ao índice de facilidade, foram usadas as seguintes cores para diferenciar o nível de dificuldade da questão :  Azul para as questões classificadas com índice muito fácil ( índice \&gt ; =0,86 ), verde para as questões classificadas com índice fácil ( 0,61 a 0,85 ), amarelo para as questões classificadas com médio ( 0,41 a 0,60 ), vermelho para as questões classificadas com difícil ( 0,16 a 0,40 ) e roxo para as questões classificadas com muito difícil ( \&lt ; =0,15 ).
Já quanto ao índice de discriminação, foram usadas as seguintes cores para qualificar a questão :  As questões classificadas com índice fraco receberam a cor vermelha ( índice \&lt ; =0,19 ), as classificadas com médio receberam a cor amarelo ( 0,20 a 0,29 ), as classificadas com bom receberam a cor verde ( 0,30 a 0,39 ) e as classificadas com muito bom ( \&gt ; =0,40 ) receberam a cor azul.
As questões objetivas do Componente de Formação Geral, segundo o índice de facilidade, foram assim avaliadas : das oito questões, uma teve o índice de facilidade classificado como muito fácil e duas como fácil.
Duas questões foram tidas como médio, por terem índice de acertos situado na faixa entre 0,41 e 0,60 ( de 41,0 \% a 60,0 \% de acertos ).
Três questões foram consideradas de dificuldade difícil, situando -se no intervalo entre 0,16 e 0,40 do índice de facilidade, ou seja, houve entre 16,0 \% e 40,0 \% de acertos.
O Gráfico 3.13, para exemplificar, analisa o comportamento da questão de número 3 de Formação Geral.
Trata -se da quinta questão mais fácil e a que obteve o maior índice de discriminação dessa parte da prova.
Neste gráfico, cada uma das cinco curvas representa o percentual de respostas em determinada alternativa da questão, em função da nota dos estudantes nesta parte da prova ( Formação Geral/Múltipla Escolha ), antes de possíveis eliminações pelo critério do ponto Como já comentado, para análise das questões objetivas relativas à Formação Geral segundo o poder de discriminação, utilizou -se, o índice de discriminação ponto bisserial.
Nesta análise, as questões foram assim avaliadas : cinco das oito questões apresentaram índices acima de 0,40 e, assim, foram classificadas com índice muito bom para esse grupo de alunos ; duas questões tiveram índice de discriminação bom, entre 0,30 e 0,39 ; e uma teve índice de discriminação médio, entre 0,20 e 0,29.
Nenhuma questão foi classificada com um índice fraco e, portanto, nenhuma questão foi eliminada pelo critério ponto bisserial.
O índice de facilidade variou de 0,19 a 0,95, e o de discriminação, de 0,26 a 0,48.
As questões com índice de discriminação muito bom, de números 3, 4, 6, 7 e 8, estão espalhadas entre os diferentes níveis de facilidade desse conjunto : duas ( questões 4 e 8 ) classificadas na categoria fácil, outras duas na categoria médio e uma ( questão 7 ) na categoria difícil do índice de facilidade.
Em particular, a questão 3 apresentou maior poder discriminatório, com índice 0,48, e foi a quinta mais fácil, com uma proporção de 0,42 de acertos.
A questão de número 5 apresentou índice de facilidade 0,19, ou seja, um quantitativo de 19,0 \% de estudantes conseguiu resolvê -la, dentro do universo de participantes.
Além disso, seu índice de discriminação foi bom, com valor igual a 0,35.
O Gráfico 3.13, para exemplificar, analisa o comportamento da questão de número 3 de Formação Geral.
Trata -se da quinta questão mais fácil e a que obteve o maior índice de discriminação dessa parte da prova.
Neste gráfico, cada uma das cinco curvas representa o percentual de respostas em determinada alternativa da questão, em função da nota dos estudantes nesta parte da prova ( Formação Geral/Múltipla Escolha ), antes de possíveis eliminações pelo critério do ponto bisserial.
A curva em vermelho corresponde à alternativa E, a correta para esta questão.
Assim, observa -se que entre os estudantes com notas mais baixas, nessa parte do exame, a situação mais frequente foi a escolha de uma das alternativas incorretas : a alternativa A ( em azul ) ou D ( em roxo ).
Na medida em que a nota aumenta, indicando desempenho melhor nesta parte da prova, aumenta concomitantemente a proporção de estudantes que selecionaram a alternativa correta E, atingindo 100 \% para os alunos com 8 acertos.
Essa análise permite verificar como a questão discriminou os grupos de desempenho, justificando o alto índice obtido na questão.
Os gráficos relativos às demais questões de Formação Geral constam do Anexo I. 3.2.2 Componente de Conhecimento Específico A Tabela 3.6 apresenta as estatísticas básicas em relação às questões objetivas do Componente de Conhecimento Específico da prova, por Grande Região.
A média do Brasil deste componente foi de 53,7.
A menor média foi observada na região Centro-Oeste ( 48,8 ) e a maior na região Sudeste ( 54,9 ).
O desvio padrão de todo o Brasil foi 15,7, sendo o menor desvio padrão encontrado na região Nordeste ( 13,7 ) e o maior na região Norte ( 16,9 ).
A mediana de todo o Brasil foi 52,6, sendo encontrada em todas as regiões do Brasil, exceto na região Sudeste ( 57,9 ).
A maior nota máxima da prova foi obtida, nas questões objetivas do Componente de Conhecimento Específico, por pelo menos um aluno na região Sudeste ( 94,7 ), enquanto a menor nota máxima foi obtida na região Nordeste ( 84,2 ).
A menor nota mínima foi obtida na região Sudeste ( 0,0 ), enquanto a maior nota mínima foi obtida nas regiões Norte e Nordeste ( 15,8 ).
A Tabela 3.7 apresenta o índice de facilidade e o índice de discriminação ( ponto bisserial ) para cada uma das questões objetivas do Componente de Conhecimento Específico.
Para facilitar a diferenciação das questões usaram -se as mesmas cores da Tabela 3.5 para as diferentes classificações dos índices de facilidade e de discriminação.
Dentre as questões objetivas da parte da prova relativa ao Componente de Conhecimento Específico, nenhuma foi anulada pela Comissão.
Desse modo, a classificação quanto ao índice de facilidade foi estabelecida com base nas 27 questões.
A partir dos índices obtidos, pode -se concluir que grande parte das 27 questões objetivas da prova foi considerada como difícil, 12 ao todo.
As 15 questões \textbf{restantes} foram classificadas quanto ao índice de facilidade da seguinte maneira : sete como fácil, outras sete como médio e uma como muito difícil.
Já quanto aos índices de discriminação das questões objetivas do Componente de Conhecimento Específico da prova, tem -se como resultado a seguinte classificação : seis das 27 questões foram consideradas como bom, enquanto quatro delas tiveram índice de discriminação muito bom.
Assim, para dez das 27 questões, os índices de discriminação foram bom ou muito bom.
Dentre as demais, nove delas foram classificadas como médio e outras oito como fraco, sendo 17, por conseguinte, a quantidade de questões nos dois patamares mais baixos de discriminação.
Constata -se, assim, que a prova – no que se refere ao Componente de Conhecimento Específico – possuía capacidade baixa de discriminar entre aqueles que dominam ou não o conteúdo.
Dentre as questões que alcançaram os maiores índices de discriminação, as de números 22, 24, 27 e 32 classificadas com índice muito bom, igual ou superior a 0,40, duas delas ( questões 22 e 27 ) foram classificadas na categoria fácil e outras duas ( questões 24 e 32 ) como médio, quanto ao índice de facilidade.
A questão de número 11 foi a mais difícil dentre as 27 questões específicas, com baixo índice de facilidade, apenas 14,0 \% de acertos.
Essa questão apresentou poder discriminatório igualmente baixo, 0,06, o que comprova ter sido esta a mais difícil para os estudantes.
Destaca -se, também, a questão 19, com índice de facilidade 0,17, o que, em termos percentuais, corresponde a 17,0 \% de estudantes que responderam acertadamente, obtendo, ainda, 0,10 de índice de discriminação.
Tais questões foram, portanto, duas das mais difíceis da prova.
As questões 11 e 19 foram eliminadas do cômputo da nota final pelo critério do ponto bisserial.
Além destas duas, as demais questões com índice fraco de discriminação, questões 9, 12, 18, 30, 31 e 33 também não foram computadas.
E outras oito como fraco, sendo 17, por conseguinte, a quantidade de questões nos dois patamares mais baixos de discriminação.
Constata -se, assim, que a prova – no que se refere ao Componente de Conhecimento Específico – possuía capacidade baixa de discriminar entre aqueles que dominam ou não o conteúdo.
Dentre as questões que alcançaram os maiores índices de discriminação, as de números 22, 24, 27 e 32 classificadas com índice muito bom, igual ou superior a 0,40, duas delas ( questões 22 e 27 ) foram classificadas na categoria fácil e outras duas ( questões 24 e 32 ) como médio, quanto ao índice de facilidade.
A questão de número 11 foi a mais difícil dentre as 27 questões específicas, com baixo índice de facilidade, apenas 14,0 \% de acertos.
Essa questão apresentou poder discriminatório igualmente baixo, 0,06, o que comprova ter sido esta a mais difícil para os estudantes.
Destaca -se, também, a questão 19, com índice de facilidade 0,17, o que, em termos percentuais, corresponde a 17,0 \% de estudantes que responderam acertadamente, obtendo, ainda, 0,10 de índice de discriminação.
Tais questões foram, portanto, duas das mais difíceis da prova.
As questões 11 e 19 foram eliminadas do cômputo da nota final pelo critério do ponto bisserial.
Além destas duas, as demais questões com índice fraco de discriminação, questões 9, 12, 18, 30, 31 e 33 também não foram computadas.
A título de exemplo das análises do comportamento das questões objetivas, o Gráfico 3.14 analisa a questão 24 do Componente de Conhecimento Específico.
Esta foi uma das questões classificadas como médio nesta prova, apresentando índice de facilidade 0,47, ou seja, 47,0 \% dos estudantes assinalaram acertadamente a opção A, correspondente ao gabarito.
Seu índice de discriminação foi igual a 0,44, classificado como muito bom.
Neste gráfico, cada uma das cinco curvas representa o percentual de respostas em determinada alternativa da questão 24, em função da nota dos estudantes nesta parte da prova, antes de possíveis eliminações de questões pelo critério do ponto bisserial.
A alternativa correta A, representada no gráfico pela curva em azul, foi escolhida em maiores proporções pelos alunos com desempenho melhor nesta parte da prova.
Já as alternativas incorretas, também denominadas distratores, foram selecionadas principalmente por aqueles com notas mais baixas.
A proporção de alunos que selecionou a resposta correta A aumenta gradativamente, chegando a atingir 100 \% para aqueles que acertaram pelo menos 21 questões, enquanto a proporção dos que escolheram alternativas incorretas decai, a partir da primeira nota não nula, como função da nota nesta parte da prova.
Os gráficos relativos às demais questões do Componente de Conhecimento Específico constam do Anexo I. 3.3 ANÁLISE DAS QUESTÕES DISCURSIVAS 3.3.1 Componente de Formação Geral As análises dos resultados de desempenho dos estudantes de Tecnologia em Agronegócio nas duas questões discursivas relativas à Formação Geral encontram -se na Tabela 3.8 e no Gráfico 3.15.
Na Tabela 3.8 observa -se que as notas médias foram mais baixas nesse conjunto de questões do que no das objetivas.
Os estudantes de todo o Brasil obtiveram, em Formação Geral, média 49,0 nas questões objetivas e 44,1 nas questões discursivas.
A mediana também confirma o pior desempenho dos alunos de todo o Brasil nas questões discursivas do Componente de Formação Geral.
Enquanto essa estatística foi de 47,5 para questões discursivas, para as questões objetivas a mediana foi de 50,0.
Pode -se, também, notar um aumento do desvio padrão de 17,9, nas questões objetivas do Componente de Formação Geral dos alunos de todo o Brasil, para 18,4 nas questões discursivas do mesmo componente.
Como já comentado, a mediana de todo o Brasil, neste componente, foi 47,5, sendo a maior mediana encontrada na região Nordeste ( 57,5 ) e a menor na região Sul ( 44,8 ).
A maior nota máxima foi obtida, nas questões discursivas do Componente de Formação Geral, por pelo menos um aluno da região Nordeste ( 94,0 ), enquanto na menor nota máxima foi obtida na região Centro-Oeste ( 74,0 ).
Em todas as regiões a nota mínima foi zero, sem exceção.
O Gráfico 3.15 representa a distribuição das notas nas questões discursivas no Componente de Formação Geral.
A moda desta distribuição ocorre no intervalo ( 50 ; 60 ].
Devido à grande quantidade de notas zero e à alta frequência de alunos que deixaram este tipo de questão em branco o intervalo [ 0 ; 10 ] se apresenta como uma segunda moda.
A distribuição possui assimetria à esquerda, coeficiente de assimetria -0,76.
Em todas as regiões o coeficiente de assimetria também é negativo, o maior em módulo, -0,90, na região Sudeste, e o menor na região Norte, -0,62.
A avaliação das questões discursivas de Formação Geral considerou em separado o conteúdo ( peso 0,8 na nota ) e o desempenho em língua portuguesa ( peso 0,2 ).
Na sequência, os resultados verificados para cada uma das questões discursivas de Formação Geral com relação ao conteúdo serão apresentados, estabelecendo -se relações com os temas abordados em cada uma delas.
Os comentários da Banca de docentes corretores a respeito do observado na correção das respostas dos estudantes, suas impressões e conclusões serão apresentados junto à análise de cada questão.
Em seguida será feita uma análise do desempenho de Língua Portuguesa.
Os comentários da Banca de docentes corretores serão apresentados para o conjunto de questões.
Cumpre esclarecer que, tendo em vista que as questões discursivas de Formação Geral são padronizadas, ou seja, constam de todas as provas, os comentários da Banca são os mesmos para todas as carreiras acadêmicas, sendo direcionados a todos os estudantes que participaram do ENADE/2013.
A seguir, serão analisados os desempenhos da Área de Tecnologia em Agronegócio nas duas questões discursivas de Formação Geral do ENADE/2013, comparando os resultados obtidos com comentários para cada questão. 3.3.1.1 Análise de Conteúdo Questão Discursiva 1 do Componente de Formação Geral Os dados de Tecnologia em Agronegócio, obtidos a partir das respostas à questão 1, encontram -se na Tabela 3.9 e no Gráfico 3.16.
Nessa questão – de melhor desempenho dentre as duas de Formação Geral – os alunos de todo Brasil tiveram média, 44,4.
A maior média para a questão 1 foi obtida na região Norte ( 54,0 ), e a menor, na região Centro-Oeste ( 41,4 ).
Quanto à variabilidade das notas, o desvio padrão de todo o Brasil foi 23,1.
O menor desvio padrão foi obtido na região Sul ( 21,9 ) e o maior desvio padrão foi obtido na região Norte ( 28,8 ).
A mediana da questão discursiva 1 foi 50,0 para as regiões Sudeste e Sul, mais alta para as regiões Norte e Nordeste ( 60,0 ) e mais baixa para a Centro-Oeste ( 45,0 ).
Para o Brasil como um todo foi também 50,0.
As notas máximas foram as mesmas para quase todas as regiões do Brasil ( 100,0 ), exceto na região Centro-Oeste ( 80,0 ).
Além disso, a nota mínima foi zero em todas as regiões do país, sem exceção.
O Gráfico 3.16 mostra a distribuição das notas na questão discursiva 1 do Componente de Formação Geral.
Observa -se que a maior frequência corresponde aos alunos que obtiveram nota no intervalo ( 40 ; 50 ], correspondendo à moda da distribuição.
Questões em branco constituem uma moda secundária.
A distribuição das notas possui assimetria negativa ( -0,46 ), assim como em todas as regiões do Brasil ( entre -0,58 para o Norte e -0,44 para o Sul ). 3.3.1.2 Comentários sobre a correção de Conteúdo das respostas à Questão Discursiva 1 O enunciado restou claro, elaborado com linguagem técnica, estando inserido nas Diretrizes gerais, com nível médio de profundidade, já que abordava tema vinculado à Área de formação da maior parte dos examinandos.
A divisão das perguntas também ficou clara indicando, passo a passo, o caminho de resposta a seguir, com base no texto.
Verificou -se dificuldade quanto à análise dos programas de governo existentes que contribuem para melhoria do saneamento básico.
Nesse ponto, os alunos referiram, com muita frequência, a ações relacionadas à Saúde, tais como : médico da família, ações educativas, ações ambientais, melhoria nas condições de habitação ( sem necessariamente vincular a programas de governo federal ou estaduais ou municipais ).
Os programas citados foram : Programa Saúde da Família ( PSF ), Programa de Aceleração do Crescimento ( PAC ) e Minha Casa Minha Vida e, ainda, o SUS.
Apesar de uma minoria de examinandos ter indicado o PAC do Governo Federal, a maior parte deles apresentou alguma solução para a questão do saneamento básico.
Os estudantes indicaram, com facilidade, a questão da necessidade do tratamento da água como item essencial do saneamento básico.
Muitos concluintes, através da sua experiência local, citaram situações em que as crianças caminhavam de pés no chão perto de rotas de esgoto a céu aberto.
Também criticaram a inércia governamental em não combater as invasões ilícitas, em locais inadequados, sem estrutura de água potável e também de saneamento básico.
Diante das observações preliminares pode ser extraída uma classificação das respostas obtidas que ficou estabelecida em fracas, médias e boas.
As notas consideradas fracas ( notas de 0 a 3,0 ) foram conferidas àqueles que não conseguiram relacionar o saneamento básico com a importância do tratamento de água, os projetos para resolver o problema e nem indicaram soluções para a questão.
As respostas que tiveram média avaliação ( notas de 3,5 a 7,0 ) tocaram, de algum modo, nos pontos \textbf{indicados}, mesmo que de forma insuficiente.
Já as respostas consideradas boas ( notas de 7,5 a 10,0 ) preencheram os requisitos do padrão de resposta, não logrando nota máxima, na maior parte dos casos, por ausência de indicação de programa governamental, mas sugerindo soluções técnicas para o problema.
As respostas mais comuns foram as consideradas medianas e que indicaram, com frequência, que o problema do saneamento básico estaria nas classes de renda mais desfavorecida e nas regiões consideradas como periferias das grandes cidades.
Também foram frequentes respostas que citaram a necessidade de tratamento de água, acompanhada da necessidade de maiores investimentos nessa área. 3.3.1.3 Análise de Conteúdo da Questão Discursiva 2 do Componente de Formação Geral A Tabela 3.10 mostra que o desempenho dos estudantes na questão 2 ( média 42,3 ) foi inferior ao obtido na questão de número 1 ( média 44,4 ).
A região Nordeste foi aquela onde a média, nessa questão, foi maior ( 55,7 ), e a de menor média foi a região Norte ( 37,7 ).
Quanto à variabilidade das notas, o desvio padrão de todo o Brasil foi 24,4, superior ao obtido na questão de número 1 ( 23,1 ).
O maior desvio nessa questão foi obtido na região Nordeste ( 28,1 ), enquanto o menor foi obtido na região Centro-Oeste ( 23,1 ).
Quanto à variabilidade das notas, o desvio padrão de todo o Brasil foi 24,4, superior ao obtido na questão de número 1 ( 23,1 ).
O maior desvio nessa questão foi obtido na região Nordeste ( 28,1 ), enquanto o menor foi obtido na região Centro-Oeste ( 23,1 ).
A mediana da questão discursiva 2 foi 45,0 para a região Sul, mais alta para as regiões Nordeste ( 60,0 ) e Sudeste ( 50,0 ) e mais baixa para as regiões Norte e Centro-Oeste ( 40,0 ).
Para o Brasil como um todo foi também 45,0.
As notas máximas ( 100,0 ) foram as mesmas em quase todas as regiões do Brasil, exceto nas regiões Norte ( 80,0 ) e Centro-Oeste ( 90,0 ).
A nota mínima foi zero em todas as regiões do Brasil, sem exceção.
O Gráfico 3.17 mostra a distribuição das notas na questão discursiva 2 do Componente de Formação Geral.
Neste gráfico, observa -se a grande quantidade de alunos que deixaram a questão 2 em branco ( em torno de 10 \% ), além da moda no intervalo ( 40 ; 50 ].
O coeficiente de assimetria das notas nesta questão é negativo ( -0,25 ), indicando que a distribuição é assimétrica à esquerda.
O mesmo ocorre para todas as regiões da Brasil ( entre -0,69 na região Nordeste e -0,22 nas regiões Sul e Centro-Oeste ). 3.3.1.4 Comentários sobre a correção de Conteúdo das respostas à Questão Discursiva 2 O tema abordado no enunciado, apesar da complexidade, é extremamente atual.
Essa atualidade facilitou a elaboração de alguma resposta, permitindo que todos os que se dispuseram a responder pudessem apresentar, pelo menos, alguns dos elementos integrantes do padrão de resposta.
A correção, seguindo o padrão de resposta, avaliou a capacidade de elaborar um texto dissertativo que abordasse o tema sob três pontos de vista : o papel da tecnologia digital na atualidade ; a garantia dos direitos do cidadão e do Estado diante do avanço tecnológico ; e o problema da segurança digital para o indivíduo, para o Estado e para as organizações ( empresas e/ou instituições nacionais ou internacionais ).
Dada a amplitude de respostas possíveis, os alunos tiveram desempenho satisfatório.
Identificamos que os concluintes, basicamente, trilharam dois caminhos : a ) a análise da influência das inovações tecnológicas, notadamente através das redes sociais, numa perspectiva essencialmente privada, daí defluindo as questões da privacidade e da segurança ; b ) uma perspectiva de segurança nacional-estatal com criticas à atuação americana de invasão de privacidade dos dados governamentais, e à inercia do governo brasileiro em reagir ao ataque, considerado, pela maioria esmagadora dos concluintes, como indevido e antiético.
Poucos foram os que uniram as duas facetas da questão : o público e o privado.
Na maior parte das provas, o concluinte indicou a ausência de segurança e de privacidade decorrentes das próprias atividades individuais.
Muitos respondentes citaram que, de forma voluntária, os próprios cidadãos enfraquecem a proteção à privacidade, pois divulgam sua vida pessoal, aspectos de seu cotidiano, por meio de imagens e do fornecimento de dados pessoais.
No aspecto segurança, foram inúmeras as indicações de acessos indevidos a dados pessoais com a perpetração de fraudes de todos os tipos, notadamente as bancárias e comerciais, com a aquisição de bens à revelia dos titulares dos dados.
Também no tema segurança, muitos indicaram o aspecto positivo decorrente da vigilância generalizada, não somente através da rede de computadores, mas também pela quantidade de câmeras existentes, que contribuem para a identificação de criminosos procurados ou dos que realizam crimes pela internet.
Na perspectiva pública, a esmagadora maioria indicou repulsa pelo comportamento americano de invasão dos dados do governo brasileiro, reclamando reação do mesmo.
Alguns poucos justificaram a intervenção americana pela “ natural ” disputa de poder pelos Estados imperialistas. 3.3.1.5 Análise de Língua Portuguesa das Questões Discursivas do Componente de Formação Geral Os dados de Tecnologia em Agronegócio, obtidos a partir das respostas às questões discursivas do Componente de Formação Geral no que tange à Língua Portuguesa, encontram -se na Tabela 3.11 e no Gráfico 3.18.
Nesse aspecto, os alunos de todo Brasil tiveram média, 47,5.
A maior média com respeito à Língua Portuguesa foi obtida na região Nordeste ( 51,7 ), e a menor, na região Centro-Oeste ( 44,3 ).
Quanto à variabilidade das notas, o desvio padrão de todo o Brasil foi 17,4.
O menor desvio padrão foi obtido na região Nordeste ( 15,2 ) e o maior desvio padrão foi obtido na região Centro-Oeste ( 18,4 ).
A mediana da nota de Língua Portuguesa foi 50,0 para as regiões Sudeste e Sul, mais alta para as regiões Norte ( 52,5 ) e Nordeste ( 55,0 ) e mais baixa para a Centro-Oeste ( 47,5 ).
Para o Brasil como um todo foi também 50,0.
A maior nota máxima foi obtida por pelo menos um aluno na região Norte ( 85,0 ), enquanto a menor nota máxima foi obtida na região Nordeste ( 75,0 ).
A nota mínima foi zero em todas as regiões do país, sem exceção.
O Gráfico 3.18 mostra a distribuição das notas de Língua Portuguesa do Componente de Formação Geral.
Observa -se que a maior frequência corresponde aos alunos que obtiveram nota no intervalo ( 50 ; 60 ], correspondendo à moda da distribuição.
Destacam -se, também o intervalo ( 40 ; 50 ] que possui distribuição muito próxima à modal, além das questões em branco que aparece como uma segunda moda.
A distribuição das notas possui assimetria negativa ( -1,22 ), assim como em todas as regiões do Brasil ( entre -1,76 para o Nordeste e -0,94 para o Norte ). 3.3.1.6 Comentários sobre a correção das respostas de Formação Geral com respeito à Língua Portuguesa Um aspecto interessante a ser destacado é a diferença de desempenho entre a questão 1 e a questão 2, principalmente nos aspectos textuais – os textos produzidos para a questão 1 apresentaram resultados melhores do que os produzidos para a questão 2.
Uma hipótese que provavelmente explique essa discrepância é o fato de o tema “ saneamento básico ” ser mais relevante para profissionais de saúde do que o tema “ segurança e privacidade ”, constituindo tema trabalhado no currículo escolar.
Prova disso é que quase todos os textos da questão 1 citam termos específicos da área de saúde, como “ contaminação feco-oral ”, “ higiene corporal ”, “ ingestão de alimentos ou água contaminada por vetores biológicos ”.
Os resultados da avaliação correspondem aos seguintes aspectos observados em cada competência : a ) aspectos ortográficos : o desempenho dos participantes ficou concentrado, majoritariamente, nos níveis 3 ( até 3 desvios ) e 2 ( de 4 a 7 desvios ), devido ao grande índice de desvios de acentuação, contrabalançada pelo pequeno número de desvios em grafia.
Foi reduzido o número de textos com ausência de desvios ( enquadrados no nível 4 ).
Alguns textos, com pior desempenho, foram enquadrados no nível 1 ( de 8 a 12 desvios ), enquanto o nível zero foi atribuído a pouquíssimos casos, já que ele revela falta absoluta de domínio das convenções ortográficas.
Observou -se, portanto, que existe uma diferença muito grande de desempenho nos dois aspectos analisados : baixíssimo índice de desvios ortográficos e grande índice de desvios de acentuação.
Em vários casos, ocorre ausência completa de acentuação gráfica.
Os resultados revelam, portanto, que a tendência dominante entre os universitários brasileiros é a eliminação da acentuação gráfica, talvez motivada pela ausência de esclarecimento dos meios de comunicação, das autoridades e das escolas sobre as decisões do Acordo Ortográfico da Língua Portuguesa de 1990.
Os casos mais sistemáticos de ausência de acentuação são :  palavras proparoxítonas ( “ politicas ”, “ publicas ”, “ juridicas ” ) ;  palavras paroxítonas terminadas em ditongo crescente ( “ varios ”, “ Brasilia ”, “ familia ” ) ;  palavras oxítonas ( “ ninguem ”, “ esta ”, “ ate ” ).
Quanto ao domínio das convenções relativas à grafia das palavras, são pontuais os problemas encontrados ( “ infecciozas ”, “ obtos ”, “ hipedemia ”, “ saniamento ”, “ cituação ”, “ abitantes ” ).
Os únicos desvios recorrentes foram : a ausência de segmentação nos casos de verbo + pronome oblíquo “ se ”, combinada, eventualmente, à alteração da grafia do pronome ( “ percebese ”, “ percebece ”, “ preferece ” ) ; a grafia da forma verbal de pretérito imperfeito do subjuntivo como se fosse uma ocorrência de verbo + pronome oblíquo ( “ distribui -se ” por “ distribuísse ” ) Vale observar, também, que, ao contrário do que se esperava, não apareceram abreviaturas próprias do “ internetês ”, ou seja, dos hábitos de comunicação escrita adquiridos pelo uso de redes sociais e e-mails. b ) aspectos textuais : o desempenho dos participantes ficou concentrado, majoritariamente, nos níveis 3 e 2, devido à grande ocorrência de problemas de estruturação textual.
Foi muito baixo o número de textos enquadrados no nível 4 por não apresentar qualquer problema estrutural.
Alguns textos com pior desempenho foram enquadrados no nível 1, enquanto o nível zero foi atribuído a pouquíssimos casos, já que ele revela a existência de um texto sem articulação e com comprometimento do sentido.
Observou -se que a grande maioria dos participantes não distribuiu as ideias em parágrafos, talvez devido ao pequeno número de linhas disponibilizadas para a resposta da questão ou, quem sabe, pela suposição de que não seria necessária essa divisão por não se tratar de um texto no modelo de uma redação dissertativo-argumentativa, como solicitado nos vestibulares.
Essa competência é a que se revela como a mais problemática entre os participantes, porque são muitos os problemas observados, desvios acumulados durante toda a formação do estudante e que não se resolvem com um estudo autodidata, como acontece com regras ortográficas ou morfossintáticas : sequência justaposta de ideias sem encaixamentos sintáticos ; redução drástica de estruturas subordinadas, ao lado do aumento na frequência de estruturas coordenadas e absolutas ; redução no uso de conectores para expressar relações lógicas essenciais à construção do texto, substituídas pela exigência de inferência por parte do interlocutor para suprir a sua ausência ; emprego equivocado de operadores que não estabelecem relações lógicas coerentes entre ideias do texto ; emprego inadequado do pronome relativo ( com omissão da preposição ou a utilização de pronome inadequado, como “ onde ” ) ; repetição exaustiva de termos sem a utilização de procedimentos mais sofisticados de substituição ( hiperonímias, hiponímias, nominalizações, expressões metafóricas ) ; frases fragmentadas que comprometem a estrutura lógico-gramatical ; frases formadas apenas por oração subordinada, sem oração principal.
Um importante aspecto a destacar é o baixíssimo desempenho de uma parte dos participantes em relação à estrutura formal do texto produzido, o que é extremamente preocupante ao se levar em conta que são graduandos em fase final de formação.
São frequentes os casos de desvios de estruturação frasal, com uso inadequado ou ausência de conectivos entre parágrafos e entre frases.
Em uma parte dos textos, falta um mínimo de textualidade e de domínio do registro padrão da língua.
Na verdade, observam -se relações linguísticas quase agramaticais, como as estabelecidas pela sequência de gerúndios sem o apoio de um ponto de partida para a organização das informações gramaticais e semânticas.
Quanto à utilização dos mecanismos de referenciação, dois aspectos devem ser destacados : a ocorrência, em uma boa parte dos textos, de repetições de palavras ou expressões sem a utilização de termos sinônimos ou pronomes, como seria adequado ; e a alternância entre os pronomes demonstrativos este / esse, e suas formas flexionadas, sem a observância de um critério de utilização.
Esse último aspecto tem se generalizado na modalidade oral da língua e, cada vez mais, tem se tornado uma constante na modalidade escrita de jornais e revistas : predominância da forma “ esse ” na modalidade oral e da forma “ este ” na modalidade escrita, provavelmente por um processo de hipercorreção. c ) aspectos morfossintáticos e vocabulares : o desempenho dos participantes ficou concentrado, majoritariamente, nos níveis 3 e 2, devido à grande ocorrência de problemas de regência, concordância e uso adequado dos sinais de pontuação.
O nível 4 foi atribuído a um número muito reduzido de textos.
Alguns textos com pior desempenho foram enquadrados no nível 1, enquanto o nível zero foi atribuído a pouquíssimos casos, já que ele revela a existência de um texto sem articulação e com comprometimento do sentido.
Os resultados são muito transparentes em relação aos aspectos mais problemáticos no desempenho dos participantes.
O desvio mais frequente é a falta do sinal indicativo da crase, o que se configura como um problema de regência, verbal ou nominal, conforme o caso – isso revela que o usuário não tem consciência de que, sob a forma do termo “ a ” existe a presença de uma preposição “ a ”, exigida pela regência do termo anterior.
Embora em outros exames, como o Enem, a falta de crase seja apenada em acentuação, nesta correção amostral esse desvio foi considerado no âmbito dos aspectos morfossintáticos.
Outro problema relacionado à regência verbal e à nominal, encontrado frequentemente nas questões, foi a ausência de preposição antes de pronome relativo, processo generalizado na modalidade oral da língua, em situações de registro informal.
Apesar da possibilidade de que essa alteração de regência se generalize no padrão escrito da Língua Portuguesa, como já está ocorrendo até em textos jornalísticos, o não emprego da preposição foi apenado neste processo de avaliação.
Outro desvio muito frequente diz respeito aos processos de concordância verbal e de concordância nominal.
Quanto à concordância de número, observou -se ausência de marca ( com sujeito anteposto ou posposto ) ou uso indevido ( uso inadequado da marca de plural comandado pelo núcleo plural da locução adjetiva, apesar de o substantivo que funciona como núcleo do sintagma nominal estar no singular ).
Uma ocorrência que se destacou foi a ausência de acento circunflexo na forma plural do presente do indicativo do verbo “ ter ”, que foi considerada como um desvio na concordância verbal e não na acentuação gráfica.
Quanto à concordância de gênero, vários casos foram observados, normalmente no âmbito de sintagmas nominais longos, em que o adjetivo está afastado do substantivo.
Quanto à questão da colocação pronominal, foram poucos os casos observados.
Apesar de serem aspectos relacionados à oralidade, concluiu -se que, no registro escrito formal, a maioria dos participantes já incorporou regras como a não introdução da frase por um pronome oblíquo e a próclise na presença de um termo atrator.
Não se adotou, entretanto, o padrão excessivamente formal descrito pelas gramáticas normativas em relação à posição do pronome oblíquo em locuções verbais, já que esse uso está muito distante da prática cotidiana, até em textos mais formais.
Quanto aos aspectos vocabulares, alguns tipos de inadequação foram observados : expressões da oralidade apareceram em algumas respostas, mas sem maior relevância do ponto de vista quantitativo ; seleção vocabular incompatível com o contexto, gerando situações de falta de inteligibilidade ; falta de domínio de vocabulário mais abstrato e de maior complexidade, essencial ao desenvolvimento do texto de base dissertativa.
Quanto à utilização dos sinais de pontuação, observou -se uma grande precariedade nos textos analisados.
Não foi apenada a ausência de vírgula para destacar locuções ou adjuntos adverbiais de pequena extensão deslocados de posição na frase, por ser um uso opcional.
São os seguintes os tipos de problemas encontrados : a ) vírgula : utilização de vírgula para separar o sujeito e o predicado ; ocorrência de apenas uma das vírgulas para separar uma palavra, uma expressão ou uma oração encaixada ; uso de vírgula no lugar do ponto para separar ideias que constituem períodos distintos ; ausência de vírgula para separar enumerações ; ausência de vírgula para separar oração adjetiva explicativa ou utilização inadequada para separar oração adjetiva restritiva ; b ) ponto e vírgula : utilização do ponto e vírgula no lugar de vírgula ; c ) ponto final : ausência de ponto final para separar períodos.
Várias marcas de oralidade foram identificadas, embora não em alta frequência : o uso do pronome relativo “ onde ” como relativo universal, falta de artigo definido antes de substantivo, repetição de palavras por falta de vocabulário, reduções como “ tá ”, “ pra ”, “ pro ”, “ prum ”, expressões informais como “ estão nem aí ”, “ tem que ”, “ levando com a barriga ”, “ não tá nem aí ”, “ a coisa fica séria ”, “ e aí haja dor de cabeça ”. 3.3.2 Componente de Conhecimento Específico Na parte da prova relativa às questões discursivas no Componente de Conhecimento Específico ( Tabela 3.12 ), observa -se que a média foi mais baixa do que nas questões discursivas do Componente de Formação Geral.
Enquanto no Componente de Formação Geral a média para estudantes de Tecnologia em Agronegócio de todo o Brasil foi 44,1, na Parte de Conhecimento Específico a média foi 31,7.
A maior média deste componente foi obtida pelos estudantes da região Sudeste ( 33,3 ), e a menor, pelos da região Norte ( 27,9 ).
Quanto à variabilidade das notas, o desvio padrão de todo o Brasil foi 15,6.
O maior desvio padrão foi encontrado na região Centro-Oeste ( 16,8 ), e o menor, na região Sul ( 14,4 ).
A maior nota máxima, nas questões discursivas do Componente de Conhecimento Específico, foi obtida na região Sudeste ( 71,7 ), enquanto a menor nota máxima foi A maior nota máxima, nas questões discursivas do Componente de Conhecimento Específico, foi obtida na região Sudeste ( 71,7 ), enquanto a menor nota máxima foi encontrada na região Norte ( 58,3 ).
Além disso, a nota mínima ( 0,0 ) foi obtida por alunos de todas as regiões do Brasil sem exceção.
A mediana do Brasil como um todo foi 33,3, também encontrada na região Centro-Oeste.
As demais medianas foram : 31,7 nas regiões Norte e Sul, 30,0 na região Nordeste e 35,0 na região Sudeste.
O Gráfico 3.19 representa a distribuição das notas nas questões discursivas no Componente de Conhecimento Específico.
Essa distribuição tem moda no intervalo de ( 30 ; 40 ].
O intervalo [ 0 ; 10 ] constitui uma moda secundária.
Nota -se que há uma tendência decrescente a partir deste intervalo até o último intervalo.
O coeficiente de assimetria é negativo para o Brasil ( -0,32 ) e para quase todas as regiões ( entre -0,42 para a região Sudeste e -0,22 para a região Centro-Oeste ), indicando uma assimetria à direita, excetuando -se apenas a região Nordeste ( 0,03 ) que possui assimetria positiva. 3.3.2.1 Análise da Questão Discursiva 3 do Componente de Conhecimento Específico Na questão 3, cujos resultados aferidos encontram -se descritos na Tabela 3.13, a média dos estudantes de todo o Brasil foi 47,1.
O desempenho médio dos estudantes do Brasil como um todo nesta questão foi superior ao desempenho médio nas questões 4 e 5.
A menor média nessa questão foi obtida pelos alunos da região Centro-Oeste ( 40,9 ), enquanto a maior média foi obtida na região Sudeste ( 50,1 ).
Quanto à variabilidade das notas, o desvio padrão de todo o Brasil foi 25,5.
O maior desvio padrão foi obtido na região Sudeste ( 26,6 ), enquanto o menor foi obtido na região Sul ( 22,5 ).
A maior nota máxima foi alcançada por pelo menos um aluno nas região Sudeste ( 100,0 ), enquanto a menor nota máxima foi obtida nas regiões Norte e Nordeste ( 90,0 ).
A mediana foi 50,0 em todas as regiões do Brasil, exceto na região Centro-Oeste ( 45,0 ).
A nota mínima do Brasil como um todo foi zero, sem exceção.
O Gráfico 3.20 mostra a distribuição das notas na questão discursiva 3, do Componente de Conhecimento Específico.
Essa distribuição é unimodal com moda no intervalo ( 40 ; 50 ].
Questões em branco constituem uma moda secundária.
O coeficiente de assimetria é negativo para todo o Brasil ( -0,38 ) e para todas as regiões ( entre -0,50 na região Sul e -0,20 na região Centro-Oeste ), indicando uma assimetria à esquerda. 3.3.2.2 Comentários sobre as respostas à Questão Discursiva 3 No enunciado da Questão 3, tanto o texto preliminar como a questão em si foram redigidos com linguagem adequada ao entendimento dos estudantes, abordando conteúdo adequado às diretrizes do curso.
O nível de conhecimento exigido era médio, já que a questão aborda assunto frequentemente tratado pela mídia, em textos técnicos e não técnicos.
O padrão de resposta era adequado, apesar de relativamente superficial em relação ao tema por não tratar da certificação de sustentabilidade, nem mencionar quem a concede, seus custos e sua relação com os sistemas de produção.
Como a questão permitia pouca flexibilidade de respostas, as diferenças entre as respostas estavam relacionadas a questões como a burocracia para obtenção das certificações, os custos, tipos de produtores nos exemplos porventura citados pelos estudantes.
Praticamente nenhum estudante respondeu de forma tão sucinta e organizada como a resposta padrão.
Apesar disso, a maioria apresentou respostas nas quais os conceitos exigidos apareciam, deixando claro que o estudante os entendia, mesmo sem enunciá -los diretamente.
Por exemplo, pouquíssimos estudantes mencionaram diretamente o conceito de agregação de valor, mas muitos mostraram por suas respostas que sabem o conceito envolvido.
Muitos estudantes apontaram a certificação como uma maneira do pequeno produtor obter maior lucratividade.
Embora isso esteja implícito na resposta padrão, já que não faz restrição de tipo de produtor, é um ponto importante a ser considerado.
Outros estudantes mencionaram as dificuldades da obtenção da certificação e do custo inerente.
A pontuação baseou -se na apresentação ( mesmo de forma não explícita ) de conceitos como : certificação, agregação de valor, nichos diferenciados de mercado, certificação como diferenciador de produtos, a disposição dos consumidores em pagar mais por produtos confiáveis.
Também foram pontuadas as respostas dos estudantes que abordaram questões como : a agricultura familiar, pequenos e médios agricultores, os custos da certificação, a burocracia envolvida e a questão da rastreabilidade de produtos.
Foram consideradas fracas as respostas que apresentavam apenas declarações relacionadas ao tema, sem responder ao que foi pedido, ou ainda apresentavam uma dissertação na qual o tema não era abordado adequadamente ( apresentam ideias soltas ou não correlacionadas ao tema ).
Por exemplo, respostas que mencionavam a importância da sustentabilidade, da preservação do ambiente etc, mas não conceituavam agregação de valor, certificação, nem mostravam como a certificação agrega valor aos produtos.
Nesses casos percebemos que o estudante não tinha conhecimento ou maturidade em relação ao tema.
As respostas medianas se caracterizaram por apresentar uma dissertação na qual o tema era abordado parcialmente ( nem todos os conceitos são apresentados ou suas relações explicadas ).
Nesses casos, percebia -se que o estudante tinha algum conhecimento sobre o assunto.
As boas dissertações abordavam adequadamente todos os aspectos pedidos - conceitos relativos à certificação, à agregação de valor e nichos específicos de mercado apresentados e suas relações explicadas.
A Questão 3 apresentou uma proporção relativamente baixa de respostas em branco, assim como desconsideradas ou nulas, o que pode ser atribuído à atualidade do tema.
Isso se reflete na presença, com muita frequência, de respostas que abordavam a busca do consumidor por uma melhor saúde através dos alimentos limpos e a preocupação com a preservação da natureza e sustentabilidade.
Os dados indicam que a maioria dos estudantes se dispôs a responder a questão e conseguiu expressar os conceitos de forma razoável, apesar de considerarmos o espaço insuficiente para uma resposta completa.
Destaca -se, ainda, que as questões objetivas 19 e 25 também podem ter contribuído para que os alunos respondessem a essa questão discursiva. 3.3.2.3 Análise da Questão Discursiva 4 do Componente de Conhecimento Específico A Tabela 3.14 contém as informações relativas à questão 4 do conjunto de questões do Componente de Conhecimento Específico.
O desempenho médio dos estudantes do Brasil como um todo nesta questão foi inferior ao desempenho na questão de número 3 e superior ao da número 5.
A média geral do Brasil foi 28,1, sendo a menor média registrada na região Sul ( 22,5 ) e a maior na região Centro-Oeste ( 32,3 ).
Quanto à variabilidade das notas, o desvio padrão de todo o Brasil foi 20,2.
O maior desvio padrão foi obtido na região Nordeste ( 24,0 ), enquanto o menor foi obtido na região Sul ( 17,2 ).
A nota máxima ( 90,0 ) foi atingida em todas as regiões, exceto na região Norte ( 80,0 ).
A mediana do Brasil como um todo foi 25,0, sendo a menor mediana encontrada nas regiões Norte e Sul ( 20,0 ) e a maior nas regiões Nordeste, Sudeste e Centro-Oeste ( 30,0 ).
A nota mínima foi zero para todas as regiões do Brasil, sem exceção.
O Gráfico 3.21 representa a distribuição de notas na questão discursiva 4, no Componente de Conhecimento Específico.
Essa distribuição é \textbf{bimodal} com a moda principal no intervalo ( 10 ; 20 ].
As questões em branco constituem uma moda secundária.
O coeficiente de assimetria é positivo para todo o Brasil ( 0,52 ) e para todas as regiões ( entre 0,18 para a região Centro-Oeste e 1,10 para a região Sul ), indicando uma assimetria à direita. 3.3.2.4 Comentários sobre as respostas à Questão Discursiva 4 Tanto o texto preliminar como a questão em si foram apresentados com linguagem adequada e clara.
O conteúdo abordado pela Questão 4 era adequado às diretrizes do curso, mas o nível de profundidade exigido foi considerado alto.
Para responder à questão eram necessários conhecimentos sobre os vários modais de transportes, suas vantagens e desvantagens e sua influência nos custos finais dos produtos.
Assuntos que, em geral, são abordados em textos técnicos, embora também apareçam na mídia.
Além disso, o enunciado permitia uma interpretação diferente da esperada, já que o texto introdutório citava composição de custos e o item a pedia a comparação e diferenciação dos modais.
Isso pode ter levado o estudante a respostas mais gerais.
A questão permitia grande flexibilidade de resposta.
No item ‘ a ’ o estudante poderia citar diferentes modais de transporte e enfatizar diferentes problemas associados aos mesmos.
No item ‘ b ’, o estudante também poderia responder de diferentes maneiras, dependendo da experiência pessoal e da região onde vive.
O item ‘ a ’ pedia a comparação e diferenciação entre os modais de transporte, e esperava -se uma resposta envolvendo todos os modais.
Pelo padrão de resposta, o estudante deveria abordar, pelo menos, dois modais de transporte e os comparar e a pontuação máxima foi obtida quando dois modais foram discutidos adequadamente.
Quase todos os estudantes apresentaram vários modais de transporte, muitos enfatizaram o modal rodoviário, que foi discutido no texto preliminar.
No item ‘ b ’, a questão era muito aberta e o que levou muitos estudantes a apresentar respostas vagas, como “ alterar as políticas do governo para o setor ” ou “ investir mais ”.
Foram raros os que apresentaram alternativas específicas, como “ construção de portos secos, próximos aos portos, para evitar as filas de caminhões ”.
Na correção do item ‘ b ’, foram exigidas, pelo menos, duas alternativas para diminuir os custos.
As respostas fracas foram aquelas nas quais o estudante apresentava apenas declarações vagas sobre o tema, ou apenas críticas gerais ao sistema de transporte.
Dissertações nas quais o tema não foi abordado adequadamente, apenas citando modais, sem dar suas características ou compará -los, também receberam pontuação baixa.
O mesmo ocorreu com respostas sem sugestões para solução dos problemas relativos ao transporte ou que o fizeram de forma muito rudimentar.
Foram consideradas medianas as respostas que abordavam o tema parcialmente, citando modais de transporte e apresentando algumas características, no entanto, com comparações fracas e sem as associar aos custos.
Com frequência as respostas foram medianas por apresentarem problemas genéricos e soluções também genéricas.
Esta questão apresentou a maior proporção de respostas em branco.
Além disso, a maioria dos estudantes que respondeu à questão, não conseguiu sintetizar o tema de forma razoável. 3.3.2.5 Análise da Questão Discursiva 5 do Componente de Conhecimento Específico A Tabela 3.15 contém as informações relativas à questão 5 do conjunto do Componente de Conhecimento Específico.
O desempenho dos estudantes nessa questão foi inferior ao das questões 3 e 4.
A nota média dos estudantes de todo o Brasil foi 19,9.
A maior média foi registrada na região Sul ( 23,5 ), enquanto a menor média foi registrada na região Nordeste ( 16,1 ).
Quanto à variabilidade das notas, o desvio padrão dos alunos do Brasil, como um todo, foi 16,2.
Enquanto o maior desvio foi encontrado na região Sul ( 17,5 ), o menor foi encontrado na região Norte ( 14,7 ).
A maior nota máxima foi alcançada nas regiões Sudeste e Sul ( 75,0 ), enquanto a menor nota máxima foi obtida na região Norte ( 55,0 ).
A mediana foi 20,0 para o Brasil como um todo, sendo a maior mediana encontrada na região Sul ( 25,0 ) e a menor na região Nordeste ( 10,0 ).
A nota mínima foi zero para todas as regiões do Brasil, sem exceção.
O Gráfico 3.22 representa a distribuição das notas da questão discursiva 5 do Componente de Conhecimento Específico.
Essa distribuição é unimodal com moda no intervalo ( 20 ; 30 ].
Além disso, as questões deixadas em branco aparecem como uma moda secundária.
O coeficiente de assimetria para todos os alunos é 0,51 e para todas as regiões se mantém positivo, variando entre 0,41 na região Sudeste e 1,01 na região Nordeste. 3.3.2.6 Comentários sobre as respostas à Questão Discursiva 5 O texto preliminar do enunciado da Questão 5 era claro, porém as perguntas ( itens ‘ a ’ e ‘ b ’ ) não deixavam claro quem deveriam ser os atores do processo.
O agricultor apenas?
Os comerciantes de agrotóxicos?
Os fabricantes?
Apesar disso, o padrão de resposta pressupunha que o ator em questão era o agricultor.
No entanto, haveria outras respostas possíveis, se, por exemplo, o estudante considerasse o revendedor, ou a própria indústria.
O assunto abordado pela questão está sujeito a leis federeis e estaduais e, portanto, pode variar de um Estado para outro.
Apesar disso, em qualquer caso, a legislação deixa claro que todos são corresponsáveis pelo destino dos resíduos produzidos, cada qual sendo responsável por ações específicas no sistema de logística reversa.
O nível de profundidade foi considerado alto, já que a resposta exigia conhecimentos sobre legislações específicas e técnicas de manuseio de produtos químicos.
Esse assunto é disciplinado pela Lei nº 7802 de 1989, alterada pela Lei nº 9.974 de 2000 e regulamentado pelo Decreto nº 4074 de 2002, capítulo IV, seção II, Art. 51 a 60, na NR-31.
Além disso, essas normas devem constar dos rótulos e bulas dos produtos ( Art. 52 e 53 do decreto nº 4074 de 2002 da Presidência da República ).
O item ‘ a ’ perguntava sobre as medidas a serem tomadas para destinar corretamente um lote de agrotóxicos fora do prazo de validade.
Esperava -se que o estudante mencionasse a necessidade de estocar em local isolado, seguro, protegido das intempéries, com piso de cimento ( isolado do solo ) longe de pessoas, animais ou alimentos e rações etc.
Deveria também lembrar a necessidade de entrar em contato com o fabricante ou distribuidor para que este providencie o recolhimento do produto. ( 2 pontos ) No item ‘ b ’ o estudante deveria explicar quais cuidados são necessários para a armazenagem e a estocagem de embalagens de agrotóxicos para sua correta destinação dentro de um sistema de logística reversa.
Como resposta esperava -se que fossem mencionadas três etapas em conformidade com a NR 31.
Foram pontuadas respostas que envolviam os seguintes cuidados : ( i ) Para embalagens rígidas contendo produtos solúveis em água : promover a tríplice lavagem, no último uso do produto, despejando a água da lavagem no tanque do pulverizador, ou a lavagem sob alta pressão ; inutilizar a embalagem furando o fundo ; não danificar os rótulos ; manter as tampas junto com as embalagens ; Para embalagens flexíveis ou contaminadas, usar sacos plásticos padronizados ; Para embalagens rígidas contendo produtos não solúveis em água e não contaminadas, manter fechadas com as próprias tampas ; Para embalagens rígidas contendo produtos não solúveis em água contaminada externamente, manter fechadas com as próprias tampas e usar sacos plásticos padronizados ; ( ii ) Armazenar em local isolado, seguro, protegido das intempéries, com piso de cimento ( isolado do solo ) longe de pessoas, animais ou alimentos e rações etc ; ( iii ) Devolver as embalagens e respectivas tampas às lojas ou postos de recebimento credenciados, no prazo de um ano, prorrogável por seis meses, caso reste produto dentro do prazo de validade ; Ou ainda, as que alertavam para : • Usar EPIs em todas as atividades com as embalagens ;  Ler o rótulo ou bula para obter as informações necessárias ; Os estudantes demostraram muita dificuldade para responder à Questão 5.
No item ‘ a ’, grande parte respondeu que o produto deveria voltar ao fabricante, mas não mencionou como armazenar o produto, nem que a responsabilidade pelo recolhimento é do fabricante ou do distribuidor.
Uns poucos basearam a resposta na destinação do produto após o recolhimento, na fábrica, citando a reciclagem de embalagens e o reaproveitamento de produtos químicos, quando possível, ou a incineração etc.
Como já mencionado a correção do item ‘ b ’ foi subdividida em três subitens.
No primeiro, muitos estudantes citaram a tríplice lavagem ou citaram o uso de sacos padronizados para embalagens contaminadas ou com produtos não solúveis em água.
No segundo, muitos citaram o armazenamento seguro, mas sem detalhar os requisitos para essa segurança.
Já para o terceiro item a ser avaliado, que a rigor não foi pedido pelo enunciado, muitos estudantes escreveram sobre a devolução das embalagens às lojas ou unidades de recebimento.
As respostas que apresentavam apenas generalidades relacionadas ao tema, mas não respondiam ao que foi pedido, receberam notas baixas.
Por exemplo, muitas vezes, o aluno citou os riscos dos agrotóxicos e a importância de devolução de embalagens com ideias soltas e sem demonstrar conhecimento legal ou técnico em relação ao tema.
Foram consideradas respostas medianas aquelas que, por exemplo, citavam que produtos vencidos deveriam ser devolvidos ao fabricante, mas não mencionavam algo sobre seu armazenamento.
Outros descreviam os procedimentos a serem tomados com as embalagens, mas deixavam escapar detalhes importantes, mostrando falta de intimidade com o tema.
Apesar de poucas, foram encontradas respostas boas, em que o estudante apresentou a forma correta de armazenar produtos vencidos, declarou que o fabricante deveria ser informado para \textbf{proceder} ao recolhimento do material e que para o item ‘ b ’ explicou os cuidados necessários para a armazenagem e a estocagem de embalagens de agrotóxicos.
Consideramos que o baixo rendimento nessa questão pode estar associado ao fato de ser a última da prova, ao mesmo tempo em que exigia uma resposta extensa para ser completa, em um espaço de apenas 15 linhas.

% matched lemmas: bimodal, indicado, proceder, restante
\end{textsample}
